\section{Discussion}
\label{sec:discussion}

The sensitivity framework of \citet{boyarsky2026} separates what pooled
experiments identify from what must be assumed. This paper adds the design question. When the answer is a wide interval,
where should the next experiment go? The question is answerable before
the experiment to a larger degree than we expected. Identification
gains are a property of the candidate's context profile relative to the
affine hull of the observed profiles. The bounds after the addition are
computable in advance up to one revealed scalar per treatment arm, with
the forecast error priced by the dual of the added constraint. Under the
maintained model a new site can only refine the set, and when a new site
instead moves conclusions the pool had pinned down, that movement is
evidence against the pooled model, concentrated exactly where the pool
claimed precision. The microcredit exercise suggests that these statements also hold in
finite samples. Computed from lending terms with the companion paper's
estimator, the predictions matched the realized experiments.

Several extensions follow the same lines. Costs enter separably, so a
budgeted rule ranks candidates by predicted gain per unit cost within
the screened pool. Choosing several sites at once is a coverage problem. The identification
gain of a set of candidates is the rank added by their profiles, which
is a matroid rank function, so greedy selection has the usual submodular
guarantee for the identification part. The width objective itself should
be re-solved at the chosen set.
When the candidate's covariate distribution is unknown, the overlap condition of \cref{subsec:augmentation} is the input most
likely to fail, and a
worst-case treatment of the candidate's covariate law would make the
precision part of the score conservative without changing the span
part. Inference for the selected design, as opposed to inference for
the bounds at a fixed design, remains open. The selection is a minimum of a maximum of estimated
program values, and the tools in the
companion paper cover each fixed program but not the selection event.

Two practical lessons follow. First, the design question depends on the choice of context basis. Every statement in this paper is
relative to the chosen context basis, and a candidate that expands the
hull in one coding can be interior in another. The companion paper's
sweeps show how much coding moves identification geometry, and a
selection exercise should be run over the same codings the final
analysis will report. Second, the most useful output of the rule can be a negative result. At Morocco the rule reports that no available candidate
delivers a signed conclusion, and that the useful experiment is one
nobody has run, in a lending environment beyond the observed hull in
Morocco's direction. A design rule that can report that no available candidate helps, and
which direction would help, gives a concrete request for the next trial.
